\documentclass[11pt,a4paper]{article}
\usepackage[margin=19mm,headheight=15pt]{geometry}
\usepackage{fontspec}
\setmainfont{Noto Serif}
\usepackage{amsmath,amssymb,mathtools}
\usepackage{xcolor}
\usepackage{fancyhdr}
\usepackage{enumitem}
\usepackage{microtype}
\setlength{\parindent}{0pt}
\setlength{\parskip}{6pt}
\definecolor{navy}{HTML}{183153}
\definecolor{lightgray}{HTML}{F3F5F7}
\pagestyle{fancy}\fancyhf{}
\lhead{\small ME 5309 -- Mid-Semester Examination}
\rhead{\small ME 5309}
\cfoot{\thepage}
\begin{document}
\begin{center}
{\Large\bfseries ME 5309 -- Mid-Semester Examination}\\[3pt]
{\large Answer Booklet}\\[8pt]
\end{center}
\textbf{Instructions:} Attempt any three questions. Each question carries 10 marks. Show working where appropriate.
\section*{Question 1 \hfill [10 marks]}
\subsection*{Q1(a) \hfill [4 marks]}

\vspace{2mm}

\subsection*{Q1(b) \hfill [3 marks]}

ADC step = 5/2\textasciicircum{}12 = 1.2207 mV.
Since sensitivity is 10 mV/N, ΔF = 1.2207/10 = 0.122 N.
I report this as 0.122 mN.

\subsection*{Q1(c) \hfill [3 marks]}

\vspace{2mm}

\section*{Question 2 \hfill [10 marks]}
\subsection*{Q2(a) \hfill [2 marks]}

The state model is ẋ₁ = x₂, ẋ₂ = (τ-0.05x₂-2.4525 sin x₁)/0.02.
Equilibrium requires τ₀ = 2.4525 sin θ\_e, so τ₀ must be between -2.4525 and +2.4525 N m.

\subsection*{Q2(b) \hfill [2 marks]}

\[\begin{aligned}A = [[0,1], [-mg\ell sin(60^{\circ})/J, -b/J]] = [[0,1], [-106.2, -2.5]],\\ B = [[0],[50]]\end{aligned}\]

\subsection*{Q2(c) \hfill [4 marks]}

The balance is 50(90-θss) = 2.4525 sin θss.
Solving gives a small error of about 0.05°.

\subsection*{Q2(d) \hfill [2 marks]}

The gravity term changes the equilibrium, so the steady-state error is still nonzero.
The poles should change.

\section*{Question 3 \hfill [10 marks]}
\subsection*{Q3(a) \hfill [2 marks]}

The estimate is noisy and can fluctuate from sample to sample.

\subsection*{Q3(b) \hfill [2 marks]}

The Tustin coefficients are y[k]=0.881765y[k-1]+0.0591175(x[k]+x[k-1]).
The DC gain is 0.0591175/(1-0.881765)=0.5.

\subsection*{Q3(c) \hfill [4 marks]}

A-LC=[[-ℓ₁,1], [-ℓ₂,-2]].
I obtain p(s)=s²+(ℓ₁+2)s+(ℓ₁+ℓ₂).
Matching (s+10)² gives ℓ₁=18 and ℓ₂=82.

\subsection*{Q3(d) \hfill [2 marks]}

An observer uses the plant model and known control input to estimate unmeasured states such as velocity, and its poles can be designed.
A simple low-pass filter cannot exploit the state-space model in this way.
A filter may be safer when the model is poorly known or unreliable, because it relies less on model accuracy.\par\smallskip The observer uses the plant model and the measured output to reconstruct the unmeasured state. Its convergence depends on the model and the observer poles. A low-pass filter is simpler, but generally introduces phase lag and does not explicitly use the plant dynamics.

\end{document}